% frontmatter/english/actividades.tex -- las actividades, una a una, para
% el adulto (frontmatter/actividades.tex, en inglés): esta página también
% tiene que caber en una (tools/check_pages.py comprueba que el día 1 está
% en la página 5).
\section*{The activities}

Each activity arrives on the day it is first used, and then comes back
from time to time. The instructions on each page, in small grey print,
can be read out by a grown-up, if needed.

{\small
\begin{multicols}{2}
\raggedright
\subsection*{From the autumn}
\begin{itemize}[leftmargin=5mm, itemsep=1pt, topsep=0pt]
\item \textbf{\lblClasifica}: read each card and match it to its box
  -- for example, \emph{I need it} or \emph{I want it}.
\item \textbf{\lblUne}: match each thing on the left to the one that
  goes with it, on the right.
\item \textbf{\lblVF}: read each sentence and circle T if it is true,
  or F if it is false.
\item \textbf{\lblTrueque}: if so many things are worth so many others,
  how many do you get for these? Drawing them helps.
\item \textbf{\lblDibuja}: draw what it says.
\item \textbf{\lblRepasa} (on Fridays): the word of the week, what they
  can now do, to tick off, and a drawing. Every ten pages there is a
  little banner to celebrate.
\item \textbf{\lblDinero}: count the notes first and then the coins,
  and write how many euros there are.
\item \textbf{\lblReparte}: share out one at a time, drawing what each
  one gets on each plate (or in each box), and see what is left over.
\item \textbf{\lblCompra}: find the price of each thing that is bought,
  and add them up.
\item \textbf{\lblProblema}: a short problem, with its sum; it can be
  drawn in the box.
\item \textbf{\lblBotes}: the money, in three jars -- to save, to spend
  and to share --; one of them is missing what is in it.
\item \textbf{\lblHucha}: save the same every week, writing down what
  there is, until the goal is reached.
\item \textbf{\lblLlega}: circle what you can buy with the money you
  have.
\item \textbf{\lblVuelta}: the change you get back is what you pay,
  take away what it costs.
\item \textbf{\lblOrdena}: put the steps in order, with a number in
  each box.
\end{itemize}
\subsection*{From the winter}
\begin{itemize}[leftmargin=5mm, itemsep=1pt, topsep=0pt]
\item \textbf{\lblElige}: there is only enough money for one thing:
  circle the one you choose, cross out the one you leave, and see what
  is left. Any choice is right.
\item \textbf{\lblCompara}: the same thing in two shops: circle the shop
  where it is cheaper, and how much you save.
\end{itemize}
\subsection*{From the spring}
\begin{itemize}[leftmargin=5mm, itemsep=1pt, topsep=0pt]
\item \textbf{\lblCuentas}: the money book: what comes in is added,
  what goes out is taken away, and in each box you write what is left.
\end{itemize}
\end{multicols}}
\newpage
